\clearpage
\section*{Appendix L\quad The Knowledge Ledgers}\label{app:ledgers}
\addcontentsline{toc}{section}{Appendix L: The Knowledge Ledgers}
\small\noindent Each claim marked \ldgmarker\ in the text rests on a field bake or the figure--theorem bake recorded in the named ledger, which ships with this work and carries the probes, the receipts, and the three registers kept apart --- what bit, what bounced, and where the boundary is. A ledger is working apparatus, not a publication, and is cited here rather than in the bibliography. This appendix is generated from the ledger index, not maintained by hand.\par\medskip
\begingroup\renewcommand{\arraystretch}{1.25}
\begin{description}
\item[\label{ldg:catastrophe_singularity}\texttt{CATASTROPHE\_SINGULARITY\_LEDGER.md}]\hfill\textsf{[field bake]}\\
The catastrophe / singularity-theory field-bake ledger --- what bit, what bounced, and the boundary. One of the three fields listed but never thrown (the overnight order named it explicitly: catastrophe $\times$54). `OWED` 622
\item[\label{ldg:variational}\texttt{VARIATIONAL\_LEDGER.md}]\hfill\textsf{[field bake]}\\
variational / action against CR --- the field the corpus uses and never names
\item[\label{ldg:index_theory}\texttt{INDEX\_THEORY\_LEDGER.md}]\hfill\textsf{[field bake]}\\
The differential-topology and index-theory field-bake ledger --- what bit, what bounced, and where the boundary is
\item[\label{ldg:integrable_systems}\texttt{INTEGRABLE\_SYSTEMS\_LEDGER.md}]\hfill\textsf{[field bake]}\\
The integrable-systems field-bake ledger --- what bit, what bounced, and where the boundary is
\end{description}
\endgroup
